\chapter{The Microscope and the Cell}\label{ch:g7:microscope}

\omterm{def:g1:living-nonliving:living}{Living things} are made of \emph{\omterm{def:g7:microscope:cell}{cells}} --- tiny units too small to see
clearly with the unaided \omterm{def:g5:nerves:senses}{eye}. The \omterm{def:g7:microscope:scope}{compound light microscope} opened this
hidden world. This chapter shows how a \omterm{prop:g4:cells:building}{microscope} works, how to use it
safely, and what \omterm{def:g1:plants-animals:animal}{animal} and \omterm{def:g1:plants-animals:plant}{plant} \omterm{def:g7:microscope:cell}{cells} look like under the lens.

\section{Seeing the invisible}

\begin{definition}[Cell]\label{def:g7:microscope:cell}
A \emph{cell}\index{cell} is the smallest unit of a living \omterm{def:g4:ecosystems:levels}{organism} that
can carry out the basic activities of life: taking in materials, using
\omterm{def:g6:nutrition:def}{energy}, growing and (usually) dividing.
\end{definition}

\begin{definition}[Microscope]\label{def:g7:microscope:scope}
A \emph{microscope}\index{microscope} is an instrument that makes small
objects appear larger and clearer. A \emph{compound light
microscope}\index{microscope!compound light} uses two lens systems
(eyepiece and objective) and \omterm{def:g4:plants-need:needs}{light} that passes through the specimen.
\end{definition}

\omimg{microscope.jpg}{0.52\linewidth}{A \omterm{def:g7:microscope:scope}{compound light microscope}: \omterm{def:g4:plants-need:needs}{light}
from below passes through the specimen and lenses to the \omterm{def:g5:nerves:senses}{eye}.}

\begin{definition}[Magnification and resolution]\label{def:g7:microscope:mag}
\emph{Magnification}\index{magnification} is how many times larger the
image is than the real object. Total magnification is approximately
\[
\text{total magnification} = \text{eyepiece} \times \text{objective}.
\]
\emph{Resolution}\index{resolution (microscope)} is the smallest distance
at which two points still look separate. High magnification is useless
without good resolution.
\end{definition}

\begin{example}[Typical powers]\label{ex:g7:microscope:powers}
Eyepiece $\times 10$ with objectives $\times 4$, $\times 10$ and
$\times 40$ gives total \omterm{def:g7:microscope:mag}{magnifications} of $40$, $100$ and $400$. Many
school \omterm{prop:g4:cells:building}{microscopes} stop near $400\times$ for clear images of \omterm{def:g7:microscope:cell}{cells}.
\end{example}

\section{Parts and light path}

\begin{method}[Using a compound microscope]\label{met:g7:microscope:use}
\begin{enumerate}
  \item Carry with two \omterm{def:g1:my-body:parts}{hands}: one on the \omterm{def:g1:my-body:parts}{arm}, one under the base.
  \item Start on the lowest-power objective; \omterm{def:g1:caring:needs}{place} the slide on the stage
        and secure it with clips.
  \item Look from the side and raise the stage (or lower the tube) until
        the objective is close but not touching the slide.
  \item Look through the eyepiece; use the coarse focus, then the fine
        focus, until the image is sharp.
  \item Centre the specimen; switch to a higher power; use \emph{only}
        fine focus at high power.
  \item Adjust the diaphragm (\omterm{def:g4:plants-need:needs}{light}) for contrast; never use coarse focus
        on high power --- you may crack the slide.
\end{enumerate}
\end{method}

\begin{omfigure}
\begin{tikzpicture}[font=\footnotesize, scale=0.95]
  \draw[thick, fill=gray!15] (0,0) rectangle (1.8,6.0);
  \draw[fill=omDef!25] (0.35,5.1) rectangle (1.45,5.7);
  \node at (0.9,5.4) {eyepiece};
  \draw[fill=omMeth!20] (0.5,3.2) rectangle (1.3,3.9);
  \node at (0.9,3.55) {objective};
  \draw[thick] (-0.2,2.5) -- (2.0,2.5);
  \draw[fill=omProp!20] (0.25,2.2) rectangle (1.55,2.5);
  \node at (0.9,1.85) {stage + slide};
  \draw[fill=omThm!25] (0.55,0.35) rectangle (1.25,1.1);
  \node at (0.9,0.7) {lamp};
  \draw[-Stealth, very thick, omExo] (2.5,0.7) -- (2.5,5.4);
  \node[omExo, align=left, anchor=west] at (2.75,3.2)
    {light path:\\lamp $\to$ specimen\\$\to$ objective\\$\to$ eyepiece\\$\to$ eye};
  \node at (0.9,-0.35) {compound microscope (schematic)};
\end{tikzpicture}

{\small \omterm{def:g4:plants-need:needs}{Light} travels upward through the specimen. The objective forms a
magnified image that the eyepiece magnifies again.}
\end{omfigure}

\section{Animal and plant cells}

\begin{definition}[Basic cell parts]\label{def:g7:microscope:parts}
Most \omterm{def:g7:microscope:cell}{cells} you will see have:
\begin{itemize}
  \item a \emph{cell membrane}\index{cell membrane} (boundary controlling
        entry and exit);
  \item \emph{cytoplasm}\index{cytoplasm} (jelly-like interior where much
        chemistry happens);
  \item a \emph{nucleus}\index{nucleus (cell)} (contains genetic
        information that \omterm{def:g6:method:control}{controls} the \omterm{def:g7:microscope:cell}{cell}).
\end{itemize}
\omterm{def:g1:plants-animals:plant}{Plant} \omterm{def:g7:microscope:cell}{cells} usually also have a rigid \emph{cell wall}\index{cell wall}
outside the membrane, a large \emph{vacuole}\index{vacuole} for storage
and support, and green \emph{chloroplasts}\index{chloroplast} that make
\omterm{def:g2:food-energy:food}{food} by \omterm{def:g5:photosynthesis:def}{photosynthesis}.
\end{definition}



\begin{omfigure}
\begin{tikzpicture}[font=\footnotesize]
  % Animal
  \draw[thick, omDef, fill=omDef!8] (-3.2,0) ellipse (2.0 and 1.5);
  \draw[thick, omThm, fill=omThm!15] (-3.5,0.2) ellipse (0.55 and 0.45);
  \node at (-3.5,0.2) {};
  \node[omThm] at (-3.5,0.95) {nucleus};
  \node at (-3.2,-1.9) {animal cell};
  \node[font=\scriptsize, align=center] at (-3.2,-2.35) {membrane\\no wall};
  % Plant
  \draw[thick, brown!60!black] (2.8,0) ellipse (2.3 and 1.7);
  \draw[thick, omDef, fill=green!10] (2.8,0) ellipse (2.1 and 1.55);
  \draw[thick, omThm, fill=omThm!15] (2.0,0.4) ellipse (0.5 and 0.4);
  \node[omThm] at (2.0,1.1) {nucleus};
  \draw[omMeth, fill=omMeth!25] (3.6,0.5) ellipse (0.45 and 0.3);
  \node[omMeth] at (3.6,1.05) {chloroplast};
  \draw[gray!60, fill=omDef!5] (3.0,-0.35) ellipse (0.85 and 0.55);
  \node at (3.0,-0.35) {vacuole};
  \node[omExo] at (4.5,1.55) {wall};
  \node at (2.8,-2.05) {plant cell};
\end{tikzpicture}

{\small Side-by-side sketch. \omterm{def:g1:plants-animals:plant}{Plant} \omterm{def:g7:microscope:cell}{cells}: wall, \omterm{def:g7:microscope:parts}{chloroplasts}, large
\omterm{def:g7:microscope:parts}{vacuole}. \omterm{def:g1:plants-animals:animal}{Animal} \omterm{def:g7:microscope:cell}{cells}: no wall or \omterm{def:g7:microscope:parts}{chloroplasts}; usually smaller \omterm{def:g7:microscope:parts}{vacuoles}.}
\end{omfigure}

\begin{proposition}[Cell theory, middle-school form]\label{prop:g7:microscope:theory}
\begin{enumerate}
  \item Living \omterm{def:g4:ecosystems:levels}{organisms} are made of one or more \omterm{def:g7:microscope:cell}{cells}.
  \item The \omterm{def:g7:microscope:cell}{cell} is the basic unit of structure and function.
  \item New \omterm{def:g7:microscope:cell}{cells} come from existing \omterm{def:g7:microscope:cell}{cells} by division.
\end{enumerate}
\end{proposition}
\admitted

\begin{example}[Onion epidermis]\label{ex:g7:microscope:onion}
A thin peel of onion \omterm{def:g5:nerves:senses}{skin} mounted in \omterm{def:g4:plants-need:needs}{water} shows brick-like \omterm{def:g1:plants-animals:plant}{plant} \omterm{def:g7:microscope:cell}{cells}
with walls and nuclei (after \omterm{def:g4:plants-need:needs}{light} staining). No green \omterm{def:g7:microscope:parts}{chloroplasts} ---
onion bulb \omterm{def:g6:uni-multi:tissue}{tissue} is not photosynthetic.
\end{example}

\section{Exercises}

\begin{exercise}[$\star$]\label{exo:g7:microscope:1}
Define \emph{\omterm{def:g7:microscope:cell}{cell}} in one sentence. Why can you not usually see individual
\omterm{def:g7:microscope:cell}{cells} without a \omterm{prop:g4:cells:building}{microscope}?
\end{exercise}

\begin{exercise}[$\star$]\label{exo:g7:microscope:2}
Name the function of: eyepiece, objective, stage, coarse focus, fine focus,
diaphragm.
\end{exercise}

\begin{exercise}[$\star$]\label{exo:g7:microscope:3}
Eyepiece $\times 10$, objective $\times 40$. What is the total
\omterm{def:g7:microscope:mag}{magnification}?
\end{exercise}

\begin{exercise}[$\star$]\label{exo:g7:microscope:4}
List three structures found in both \omterm{def:g1:plants-animals:animal}{animal} and \omterm{def:g1:plants-animals:plant}{plant} \omterm{def:g7:microscope:cell}{cells}. List three
structures typical of \omterm{def:g1:plants-animals:plant}{plant} \omterm{def:g7:microscope:cell}{cells} but not of \omterm{def:g1:plants-animals:animal}{animal} \omterm{def:g7:microscope:cell}{cells}.
\end{exercise}

\begin{exercise}[$\star$]\label{exo:g7:microscope:5}
State the three ideas of \omterm{def:g7:microscope:cell}{cell} theory (\cref{prop:g7:microscope:theory}).
\end{exercise}

